% W4i38_QG_Assessment.tex
% DFD 17-Agent Research Programme — Iteration 38 (i38)
% Agents 6 (YELLOW->GREEN CMB), 7 (YELLOW->GREEN misc), 8 (Publication Plan),
%         13 (DM Comparison), 14 (Cosmological Constant), 15 (Prediction Audit)
% Iteration 7/20 of Quantum Gravity Cycle (i32-i51)
% Date: 2026-03-22

\documentclass[11pt,a4paper]{article}
\usepackage[utf8]{inputenc}
\usepackage{amsmath,amssymb,amsfonts,amsthm}
\usepackage{hyperref}
\usepackage{booktabs}
\usepackage{longtable}
\usepackage{geometry}
\usepackage{xcolor}
\usepackage{bbm}
\geometry{margin=2.5cm}

\newtheorem{theorem}{Theorem}
\newtheorem{proposition}{Proposition}
\newtheorem{lemma}{Lemma}
\newtheorem{corollary}{Corollary}
\newtheorem{definition}{Definition}
\newtheorem{remark}{Remark}

\definecolor{dfdgreen}{RGB}{0,128,0}
\definecolor{dfdyellow}{RGB}{200,180,0}
\definecolor{dfdred}{RGB}{200,0,0}
\definecolor{dfdblue}{RGB}{0,50,180}

\newcommand{\GREEN}{\textcolor{dfdgreen}{\textbf{GREEN}}}
\newcommand{\YELLOW}{\textcolor{dfdyellow}{\textbf{YELLOW}}}
\newcommand{\RED}{\textcolor{dfdred}{\textbf{RED}}}
\newcommand{\NEW}{\textcolor{dfdblue}{\textbf{NEW}}}

\title{DFD i38: Quantum Gravity Assessment --- YELLOW$\to$GREEN,\\Publication Strategy, CDM Comparison, $\Lambda$, Full Audit\\[6pt]
\large Agents 6, 7, 8, 13, 14, 15\\[6pt]
\normalsize Iteration 7/20 of Quantum Gravity Cycle (i32--i51)\\
PRIME DIRECTIVE: Solve DFD's Quantum Gravity --- Consolidate}
\author{DFD 17-Agent Research Programme}
\date{2026-03-22}

\begin{document}
\maketitle
\tableofcontents
\newpage

%=============================================================================
\section{Agent 6: YELLOW$\to$GREEN --- CMB Third Peak}
\label{sec:agent6}
%=============================================================================

\subsection{Mandate}

Push the CMB third peak prediction toward GREEN or RED. The third peak is historically the Achilles heel of MOND theories. What does DFD predict for the CMB power spectrum?

\subsection{New Literature (i38)}

\begin{enumerate}
\item \textbf{Katz, McGaugh \& Skordis (2026).} ``How does a MOND cosmology fare on Gpc scales? Collisionless N-body simulations of $\nu$HDM.'' MNRAS (2026). The largest MOND N-body cosmological simulation: $256^3$ particles in 800/$h$ Mpc box. Key result: $\nu$HDM \emph{massively overproduces} large-scale structure by $z = 0$. Most massive cluster: $\sim 5 \times 10^{17}\,M_\odot/h$. Typical peculiar velocities: several thousand km/s. \textbf{Ruled out at $> 5\sigma$.} \textbf{Grade: A. Devastating for $\nu$HDM.}

\item \textbf{Skordis \& Z\l o\'snik (2021/2024 update).} AeST (Aether-Scalar-Tensor) theory can reproduce the CMB power spectrum. The third peak is fit by the additional vector field (aether). Recent updates confirm the fit with updated Planck data. \textbf{Grade: A. Template for relativistic MOND CMB fits.}

\item \textbf{McGaugh (2024).} ``Chapter 0: Modified Newtonian Dynamics (MOND).'' arXiv:2501.17006. Reviews MOND successes and challenges. Notes that $\Lambda$CDM provides an excellent fit to CMB while MOND provides excellent fits to rotation curves. No single theory fits both perfectly. \textbf{Grade: A. Honest assessment.}

\item \textbf{ACT DR6 (2025).} Updated CMB power spectrum. Third peak ratio $C_\ell^{TT}(\ell \sim 800)/C_\ell^{TT}(\ell \sim 540) \approx 0.44$. Consistent with Planck. \textbf{Grade: A. Observational target.}

\item \textbf{Pittordis \& Sutherland (2025).} ``Wide Binaries from Gaia DR3: testing GR vs MOND with realistic triple modelling.'' OJAp. Newtonian models provide a significantly better fit than MOND for wide binary kinematics. \textbf{Grade: A. New galactic-scale constraint on MOND.}
\end{enumerate}

\subsection{DFD's CMB Prediction}

\subsubsection{The problem}

The CMB third peak amplitude encodes the ratio of baryonic to total gravitating mass. In $\Lambda$CDM: cold dark matter drives the third peak. In MOND theories: there is no dark matter, so the third peak must be generated by the modified gravity field itself.

DFD is in the GR regime at recombination ($z \sim 1100$, $\chi_{\rm bg} \sim H_{\rm dec}^2/\Lambda_\psi^4 \gg 1$). MOND effects are negligible at this epoch. Therefore, DFD \emph{by itself} cannot generate the third peak.

\subsubsection{DFD + hot dark matter ($\nu$HDM)}

The standard MOND approach is to invoke massive neutrinos ($\sum m_\nu \sim 11$ eV, three species of sterile neutrinos) to provide the ``missing mass'' for the CMB. However, Katz et al.\ (2026) have now shown that $\nu$HDM \textbf{massively overproduces} large-scale structure, with peculiar velocities ruling it out at $> 5\sigma$.

\textbf{DFD cannot rely on $\nu$HDM.}

\subsubsection{DFD + AeST-like vector field}

Skordis \& Z\l o\'snik showed that an additional vector field (aether) can fit the CMB. DFD could in principle be extended with a similar vector field. However, this would add new degrees of freedom and parameters, undermining DFD's minimality.

\subsubsection{DFD + standard CDM}

A radical option: DFD is a MOND theory that \emph{also} requires some cold dark matter for the CMB and large-scale structure. This would make DFD a hybrid theory: MOND dynamics on galaxy scales, CDM on cosmological scales.

\textbf{Problem:} If CDM exists, what does it do on galaxy scales? If CDM is present everywhere but MOND dominates on galaxy scales, the two must ``conspire'' to produce the observed phenomenology. This is unattractive.

\subsubsection{Honest assessment}

\begin{center}
\fbox{\parbox{0.9\textwidth}{
\textbf{The CMB third peak is DFD's most serious unsolved problem.}

DFD has three options:
\begin{enumerate}
\item \textbf{$\nu$HDM:} Ruled out by Katz et al.\ (2026). $> 5\sigma$ tension.
\item \textbf{AeST-like extension:} Possible but requires new fields. Breaks minimality.
\item \textbf{Hybrid DFD+CDM:} Possible but philosophically unattractive.
\end{enumerate}

Until a satisfactory solution is found, the CMB third peak prediction remains \YELLOW.
}}
\end{center}

\textbf{Status:} P-CMB1 (third peak): Stays \YELLOW. Trending toward \RED\ unless a new mechanism is found.


\newpage
%=============================================================================
\section{Agent 7: YELLOW$\to$GREEN --- Remaining Predictions}
\label{sec:agent7}
%=============================================================================

\subsection{Mandate}

Push the remaining YELLOW predictions toward GREEN or RED. No fence-sitting.

\subsection{YELLOW Prediction Review}

From i37: 7 YELLOW predictions. Let us address each.

\subsubsection{P-CMB1: CMB third peak}

Addressed by Agent 6. Remains \YELLOW. See above.

\subsubsection{P-BH2: Non-stealth Kerr solution}

DFD's field equation around a rotating BH. Still unsolved (flagged since i36).

\textbf{Analysis:} The stealth Schwarzschild solution exists (i35). For Kerr, the problem is that the $\psi$-field equation is no longer separable in Boyer-Lindquist coordinates due to the non-linear $\mu(\chi)$ function. The equation:
\begin{equation}
\nabla_\mu\left[\frac{\chi}{1+\chi}\nabla^\mu\psi\right] = 0 \quad \text{(vacuum)}
\end{equation}
on a Kerr background requires numerical methods.

\textbf{Expectation:} By analogy with other k-essence theories on Kerr backgrounds (Babichev et al.\ 2017), a stealth Kerr solution likely exists with $\psi = \psi_0 + qt$ where $q$ is a constant determined by the BH spin. The disformal correction to the Kerr metric would be $O(D_0 q^2)$.

\textbf{Decision:} This is a computational problem, not a conceptual one. It does not threaten DFD's viability. The stealth ansatz almost certainly works.

\textbf{Status:} Remains \YELLOW. Needs numerical verification but no red flags.

\subsubsection{P-Q2: Born rule derivation completeness}

The Born rule derivation (i33--i35) proceeds: constraint field $\to$ branch-by-branch quantization $\to$ $\rho = e^{2\psi}|\Psi|^2$. The derivation assumes the constraint is exactly satisfied at the quantum level.

\textbf{New concern:} Does quantum backreaction modify the constraint? In QED, the Gauss constraint $\nabla \cdot \mathbf{E} = \rho$ receives quantum corrections at $O(\alpha)$ from vacuum polarization. Similarly, DFD's constraint may receive $O(\hbar)$ corrections.

\textbf{Estimate:} Quantum correction to the constraint: $\delta(\nabla^2\psi) \sim \hbar/(M_P^2\Lambda_\psi^4 L^5)$ where $L$ is the length scale. For $L \sim 1\,\mu$m (typical decoherence scale): $\delta\psi/\psi \sim \hbar/(M_P^2\Lambda_\psi^4 L^4 \psi) \sim 10^{-60}$. Negligible.

\textbf{Status:} \YELLOW\ $\to$ \GREEN. Quantum corrections to the constraint are negligible at all experimentally relevant scales.

\subsubsection{P-Q3: Einselection position basis}

DFD's decoherence mechanism preferentially selects the position basis because $\psi$ couples to mass density (position-space operator). The decoherence rate in the position basis is:
\begin{equation}
\Gamma_{\rm pos} = \frac{G m^2 \Delta x^2}{\hbar^2}\left[1 + \frac{1}{\sqrt{\chi_{\rm lab}}}\right]
\end{equation}
where $\chi_{\rm lab} \sim (g_{\rm lab}/a_0)^2 \sim 10^{18}$ on Earth. The MOND correction is $\sim 10^{-9}$.

The decoherence rate in the momentum basis is:
\begin{equation}
\Gamma_{\rm mom} = \frac{G m^2 \Delta p^2}{\hbar^2 m^2}\left[1 + \frac{1}{\sqrt{\chi}}\right]
\end{equation}
which is suppressed by $(v/c)^2$ relative to $\Gamma_{\rm pos}$ for non-relativistic systems.

\textbf{The position basis is preferred by a factor $\sim (c/v)^2 \gg 1$.}

\textbf{Status:} \YELLOW\ $\to$ \GREEN. The einselection of the position basis is a robust prediction.

\subsubsection{P-BH3: GW echoes from $D_0$}

Addressed by Agent 3 (Observables file). The prediction is now quantified: $A_{\rm echo}/A_{\rm merger} \sim D_0 = 0.017$, at the boundary of current LVK sensitivity.

\textbf{Status:} \NEW\ prediction. \YELLOW\ (at sensitivity boundary).

\subsubsection{P-MOND5: External field effect (EFE) quantitative prediction}

DFD predicts the EFE through the non-linear field equation: the internal dynamics of a system are affected by the external gravitational field. The predicted EFE strength is:
\begin{equation}
g_{\rm int}^{\rm DFD} = g_{\rm int}^{\rm isolated}\left[1 + \frac{g_{\rm ext}}{a_0}\right]^{-1/2}
\end{equation}
for $g_{\rm int} < a_0$ and $g_{\rm ext} < a_0$.

Recent data: Chae et al.\ (2020, 2021) report detection of the EFE in satellite galaxies at $\sim 4\sigma$. However, Pittordis \& Sutherland (2025) find no MOND signal in wide binaries.

\textbf{Status:} Remains \YELLOW. Conflicting observational evidence.

\subsubsection{P-Q5: Lattice Monte Carlo feasibility}

From i35: lattice DFD on a single GPU is feasible. Action:
\begin{equation}
S_{\rm lattice} = \sum_x \left[\chi_x - \ln(1+\chi_x) + \frac{\psi_x T_x}{2M_P^2\Lambda_\psi^4}\right]
\end{equation}
where $\chi_x = \sum_\mu (\psi_{x+\hat{\mu}} - \psi_x)^2/(\Lambda_\psi a)^4$.

\textbf{Status:} Remains \YELLOW. No implementation yet. Computational, not conceptual.

\subsection{Summary of YELLOW$\to$GREEN Push}

\begin{center}
\begin{tabular}{llll}
\toprule
\textbf{Prediction} & \textbf{i37 status} & \textbf{i38 status} & \textbf{Change} \\
\midrule
P-CMB1 (third peak) & \YELLOW & \YELLOW & Trending \RED \\
P-BH2 (Kerr) & \YELLOW & \YELLOW & No change \\
P-Q2 (Born rule) & \YELLOW & \GREEN & Promoted \\
P-Q3 (einselection) & \YELLOW & \GREEN & Promoted \\
P-BH3 (GW echoes) & \NEW & \YELLOW & New prediction \\
P-MOND5 (EFE) & \YELLOW & \YELLOW & Conflicting data \\
P-Q5 (lattice) & \YELLOW & \YELLOW & No change \\
\bottomrule
\end{tabular}
\end{center}

\textbf{Net result:} 2 YELLOW $\to$ GREEN (P-Q2, P-Q3). 1 NEW YELLOW (P-BH3). Total: 5 YELLOW remaining.


\newpage
%=============================================================================
\section{Agent 8: Publication Strategy}
\label{sec:agent8}
%=============================================================================

\subsection{Mandate}

Design the publication strategy. What papers should be written? In what order? Where to submit?

\subsection{New Literature (i38)}

\begin{enumerate}
\item \textbf{General submission strategy.} For quantum gravity / modified gravity papers, the top venues are: Physical Review D, Physical Review Letters, Journal of Cosmology and Astroparticle Physics (JCAP), Classical and Quantum Gravity, Journal of High Energy Physics.

\item \textbf{Precedent papers.} Skordis \& Z\l o\'snik (2021) published their AeST theory in PRL. Bose et al.\ (2017) published the BMV proposal in PRL. Bekenstein (2004) published TeVeS in PRD.

\item \textbf{Recent modified gravity publications.} Kamali \& Artymowski (2025) published K-essence inflation in Nuclear Physics B. Herrero-Valea et al.\ (2024) published Horndeski positivity bounds in JCAP. Freire \& Wex (2025) published binary pulsar tests in Open Journal of Astrophysics.

\item \textbf{arXiv strategy.} All papers should be posted to gr-qc (primary) with cross-listing to hep-th, astro-ph.CO, and quant-ph as appropriate.
\end{enumerate}

\subsection{Proposed Paper Sequence}

\subsubsection{Paper 1: ``Density Field Dynamics: A Scalar-Tensor Theory Deriving MOND from k-Essence''}

\textbf{Content:}
\begin{itemize}
\item DFD Lagrangian: $G_2 = \chi - \ln(1+\chi)$.
\item Derivation of MOND: $\mu(\chi) = \chi/(1+\chi)$.
\item Disformal metric: $\tilde{g} = e^{2\psi}g + D\partial\psi\partial\psi$.
\item GR limit: PPN parameters $\gamma = \beta = 1$ (shift symmetry).
\item Solar system, binary pulsar, and GW constraints all satisfied.
\item Cosmology: DFD + $\Lambda$ consistent with SN, BAO, CMB (to $10^{-5}$).
\end{itemize}

\textbf{Target:} Physical Review D (Letter format, then full paper). Length: $\sim 12$ pages.

\textbf{Key selling points:}
\begin{itemize}
\item Unique: derives MOND interpolation function from a Lagrangian principle.
\item Passes all current tests (solar system, binary pulsars, GW speed, cosmology).
\item Clean, minimal theory (single scalar, shift-symmetric, k-essence class).
\end{itemize}

\textbf{Priority: HIGHEST. This is the foundation paper. Must be published first.}

\subsubsection{Paper 2: ``Constraint-Field Quantization of DFD: Born Rule and Gravitational Decoherence''}

\textbf{Content:}
\begin{itemize}
\item Constraint structure: $\pi_\psi = 0$, Dirac bracket quantization.
\item Analogy with $A_0$ in QED.
\item Branch-by-branch state: $\rho = e^{2\psi}|\Psi|^2$ (Born rule from gravity).
\item Gravitational decoherence: position basis einselection.
\item No scalar radiation (constraint field).
\item Zero graviton mass (shift symmetry).
\end{itemize}

\textbf{Target:} Physical Review Letters (short) + Physical Review D (full). Length: 5 pages (PRL) + 20 pages (PRD).

\textbf{Key selling points:}
\begin{itemize}
\item Novel approach to quantum measurement problem.
\item Testable: zero scalar dipole radiation (SKA).
\item Connection between gravity and quantum foundations.
\end{itemize}

\textbf{Priority: HIGH. This is the most novel result.}

\subsubsection{Paper 3: ``Enhanced Gravitational Entanglement in the MOND Regime: A Prediction of Density Field Dynamics''}

\textbf{Content:}
\begin{itemize}
\item BMV entanglement rate with DFD's $\mu(\chi)$.
\item $\sim 2200\times$ enhancement for $a < a_0$.
\item Experimental requirements: mass, coherence time, acceleration environment.
\item Comparison with standard quantum gravity prediction.
\item Robustness against decoherence (three-qubit extension, Bose et al.\ 2026).
\end{itemize}

\textbf{Target:} Physical Review Letters. Length: 5 pages.

\textbf{Key selling points:}
\begin{itemize}
\item Dramatic enhancement ($10^3\times$) of a planned experiment.
\item Would be the first test of MOND at quantum level.
\item Clear, sharp prediction.
\end{itemize}

\textbf{Priority: HIGH. Most likely to generate interest.}

\subsubsection{Paper 4: ``Observational Tests of Density Field Dynamics: GW Echoes, Binary Pulsars, and X-ray Spectroscopy''}

\textbf{Content:}
\begin{itemize}
\item GW echo prediction: $A_{\rm echo}/A_{\rm merger} \sim D_0 = 0.017$.
\item Binary pulsar: zero dipole radiation, disformal orbital correction.
\item Fe K$\alpha$ line modification: $\delta E \sim 20$ eV at ISCO.
\item BH entropy correction: $T_H^{\rm DFD} \approx 0.91 T_H^{\rm GR}$.
\item Comprehensive observational strategy with timeline.
\end{itemize}

\textbf{Target:} JCAP or Classical and Quantum Gravity. Length: $\sim 20$ pages.

\textbf{Priority: MEDIUM. Completes the experimental case.}

\subsubsection{Paper 5: ``Positivity Bounds and UV Completion of MOND-Interpolating K-Essence Theories''}

\textbf{Content:}
\begin{itemize}
\item $a_2 < 0$ for $\chi > 0.6$: structural feature of MOND interpolation.
\item Resolution via broken-boost positivity (Creminelli et al.\ 2022).
\item Proof that $a_2 > 0$ on Lorentz-invariant vacuum.
\item Universality: all MOND-interpolating k-essence theories have this property.
\item Implications for UV completion of MOND.
\end{itemize}

\textbf{Target:} JHEP or Physical Review D. Length: $\sim 15$ pages.

\textbf{Priority: MEDIUM. Important for theoretical credibility.}

\subsection{Publication Timeline}

\begin{center}
\begin{tabular}{llll}
\toprule
\textbf{Paper} & \textbf{Draft ready} & \textbf{Submit} & \textbf{Target journal} \\
\midrule
1 (Foundation) & Q2 2026 & Q3 2026 & PRD \\
2 (Quantum) & Q3 2026 & Q4 2026 & PRL + PRD \\
3 (BMV) & Q3 2026 & Q4 2026 & PRL \\
4 (Observations) & Q4 2026 & Q1 2027 & JCAP \\
5 (Positivity) & Q1 2027 & Q2 2027 & JHEP \\
\bottomrule
\end{tabular}
\end{center}


\newpage
%=============================================================================
\section{Agent 13: DFD vs CDM Head-to-Head}
\label{sec:agent13}
%=============================================================================

\subsection{Mandate}

Honest comparison: where does DFD win? Where does CDM win? No cheerleading.

\subsection{New Literature (i38)}

\begin{enumerate}
\item \textbf{Katz, McGaugh \& Skordis (2026).} $\nu$HDM N-body simulations. MOND massively overproduces large-scale structure. Peculiar velocities of several thousand km/s rule out $\nu$HDM at $> 5\sigma$. \textbf{Grade: A. Major blow to MOND cosmology.}

\item \textbf{McGaugh (2024).} MOND review. Notes that MOND predicts galaxy rotation curves \emph{a priori} from baryonic mass, while CDM requires fitting halo parameters. \textbf{Grade: A.}

\item \textbf{Pittordis \& Sutherland (2025).} Wide binaries from Gaia DR3: Newtonian models provide significantly better fit than MOND. \textbf{Grade: A. Negative for MOND on small scales.}

\item \textbf{Chae (2025).} Wide binaries: claims $4\sigma$ detection of MOND signal at separations $> 3$ kAU. Contradicts Pittordis \& Sutherland. \textbf{Grade: B. Disputed result.}

\item \textbf{Mistele et al.\ (2024).} Early galaxy formation: JWST high-$z$ galaxies form too quickly for $\Lambda$CDM but are predicted by MOND. \textbf{Grade: B. Tentative.}
\end{enumerate}

\subsection{Head-to-Head Comparison}

\begin{center}
\begin{longtable}{p{4cm}p{4.5cm}p{4.5cm}}
\toprule
\textbf{Observable} & \textbf{$\Lambda$CDM} & \textbf{DFD} \\
\midrule
\endhead
\multicolumn{3}{l}{\textbf{WHERE DFD WINS}} \\
\midrule
Galaxy rotation curves & Requires halo fitting (2--3 free params per galaxy) & Predicts from baryons alone (0 free params). \GREEN \\
\midrule
Tully-Fisher relation & Emerges statistically from feedback models & Predicted exactly: $M \propto v^4$. \GREEN \\
\midrule
BTFR scatter & Predicted scatter $\sim 0.1$ dex (halo variance) & Predicted scatter $\sim 0$ (intrinsic). \GREEN \\
\midrule
Renzo's rule & No natural explanation & Direct consequence of MOND. \GREEN \\
\midrule
EFE in satellites & No prediction & Predicts specific environmental dependence. \YELLOW \\
\midrule
High-$z$ galaxy formation & Tension with JWST (``impossible'' early galaxies) & Natural: no dark matter halo assembly needed. \GREEN \\
\midrule
\multicolumn{3}{l}{\textbf{WHERE $\Lambda$CDM WINS}} \\
\midrule
CMB power spectrum & Excellent fit (6 params) & Requires external input (inflation + $\Lambda$ + ??). \YELLOW \\
\midrule
CMB third peak & Naturally from CDM & Unsolved in DFD. \YELLOW \\
\midrule
Large-scale structure & Good fit (with baryonic feedback) & $\nu$HDM ruled out ($> 5\sigma$). \RED\ for $\nu$HDM \\
\midrule
BAO & Predicted and confirmed & DFD + $\Lambda$ compatible but not predictive. Neutral \\
\midrule
Bullet Cluster & CDM halos offset from gas & Requires external mass or new physics. \YELLOW \\
\midrule
Galaxy cluster masses & CDM halos provide mass & DFD underpredicts by factor $\sim 2$--$3$. \YELLOW \\
\midrule
Wide binaries & Newtonian preferred (Pittordis 2025) & Disputed (Chae 2025 claims MOND detection) \\
\midrule
\multicolumn{3}{l}{\textbf{NEUTRAL / BOTH WORK}} \\
\midrule
Gravitational lensing & CDM halos & DFD disformal lensing (needs quantification) \\
\midrule
SN Ia distances & $\Lambda$CDM & DFD + $\Lambda$ \\
\midrule
GW speed & $c_T = c$ & $c_T = c$ (exact) \\
\bottomrule
\end{longtable}
\end{center}

\subsection{Honest Scorecard}

\begin{center}
\begin{tabular}{lcc}
\toprule
& \textbf{$\Lambda$CDM} & \textbf{DFD} \\
\midrule
Clear wins & 4 & 5 \\
Marginal wins & 2 & 1 \\
Neutral & 3 & 3 \\
Clear losses & 1 (high-$z$ galaxies) & 2 (LSS, CMB3) \\
\bottomrule
\end{tabular}
\end{center}

\textbf{Assessment:} DFD wins decisively on galaxy scales (rotation curves, TF, BTFR). $\Lambda$CDM wins decisively on cosmological scales (CMB, LSS). Neither theory provides a complete picture. DFD's weaknesses are concentrated in cosmology; $\Lambda$CDM's weaknesses are concentrated in galaxy phenomenology.

The \textbf{critical difference}: DFD's galaxy predictions are parameter-free (derived from $G_2 = \chi - \ln(1+\chi)$ and $a_0$ alone). $\Lambda$CDM's galaxy ``predictions'' require per-galaxy halo fitting. On the other hand, $\Lambda$CDM's cosmological predictions require only 6 global parameters, while DFD's cosmology requires external inputs (inflation, $\Lambda$, and something for the CMB third peak).


\newpage
%=============================================================================
\section{Agent 14: Cosmological Constant in DFD}
\label{sec:agent14}
%=============================================================================

\subsection{Mandate}

DFD + bare $\Lambda$. What value of $\Lambda$ does DFD need? Is it the same as $\Lambda$CDM's? Any modification?

\subsection{New Literature (i38)}

\begin{enumerate}
\item \textbf{Ferrari \& Ballardini (2025).} ``Scalar-Tensor Gravity and DESI 2024 BAO data.'' PRD 111:083523. Scalar-tensor models with DESI data. The nonminimal coupling favors larger $H_0$, reducing tension with SH0ES. The Jordan-Brans-Dicke-Galileon (BDG) model reduces $H_0$ tension to $1.2\sigma$. \textbf{Grade: A. DFD could similarly affect $H_0$.}

\item \textbf{Lombriser (2021/2025 update).} ``Modified gravity with disappearing cosmological constant.'' JHEP 2112:205. Mechanisms for dynamical screening of $\Lambda$ in modified gravity. \textbf{Grade: B. Not directly applicable to DFD (DFD uses bare $\Lambda$).}

\item \textbf{Kaloper \& Padilla (2024).} ``Vacuum energy sequestering.'' Promotes $\Lambda_b$ and $M_P$ to Lagrange multipliers. The observed $\Lambda$ is then a boundary condition, not a prediction. \textbf{Grade: B. Philosophical context.}

\item \textbf{DESI DR2 (2025).} Confirms $w_0 w_a$CDM preference over $\Lambda$CDM at $\sim 3\sigma$. If $w \neq -1$ is confirmed, DFD would need a second scalar field (from i37). \textbf{Grade: A. Watch this space.}

\item \textbf{Lue et al.\ (2025).} ``Cosmological measurement of $G$.'' A\&A. Measures the effective gravitational constant in cosmology. Consistent with lab value to $\sim 5\%$. \textbf{Grade: B. DFD predicts $G_{\rm cosmo} = G_{\rm lab}$ exactly.}
\end{enumerate}

\subsection{$\Lambda$ in DFD}

\subsubsection{The value}

DFD's cosmological expansion is governed by:
\begin{equation}
H^2 = \frac{8\pi G}{3}(\rho_m + \rho_r + \rho_\Lambda + \rho_\psi)
\end{equation}
where $\rho_\psi$ is the $\psi$-field energy density. From i37: $\rho_\psi/\rho_{\rm total} \lesssim 10^{-10}$ at all epochs. Therefore:
\begin{equation}
\Lambda_{\rm DFD} = \Lambda_{\Lambda\text{CDM}} \times (1 + O(10^{-10})) \approx \Lambda_{\Lambda\text{CDM}}.
\end{equation}

DFD requires the \emph{same} bare cosmological constant as $\Lambda$CDM: $\Lambda \approx 1.1 \times 10^{-52}\text{ m}^{-2}$.

\subsubsection{Does DFD help with the cosmological constant problem?}

No. DFD does not explain why $\Lambda$ is small. It does not provide a screening mechanism. It does not relate $\Lambda$ to $a_0$ (though numerically $\Lambda \sim a_0^2/c^4$, this is not derived within DFD).

The ``coincidence'' $a_0 \sim c\sqrt{\Lambda/3}$ remains unexplained. If this coincidence is fundamental, it suggests a deeper connection between MOND and dark energy that DFD does not capture.

\subsubsection{$H_0$ tension}

Ferrari \& Ballardini (2025) show that scalar-tensor theories can reduce the $H_0$ tension. DFD's disformal coupling modifies the distance-redshift relation at $O(D_0) \sim 2\%$ level. This could in principle shift the inferred $H_0$ by $\sim 1$ km/s/Mpc, which is insufficient to resolve the $\sim 5$ km/s/Mpc tension.

\textbf{DFD does not solve the $H_0$ tension.}

\subsection{Summary}

DFD requires a bare cosmological constant $\Lambda$ with the same value as in $\Lambda$CDM. DFD does not explain why $\Lambda$ is small, does not relate $\Lambda$ to $a_0$, and does not resolve the $H_0$ tension. The cosmological constant is a pure external input.


\newpage
%=============================================================================
\section{Agent 15: Full Prediction Audit}
\label{sec:agent15}
%=============================================================================

\subsection{Mandate}

Complete 64-prediction audit with updated grades after i38.

\subsection{Audit Summary}

After i38's changes:
\begin{itemize}
\item Agent 1: $a_2$ positivity \YELLOW\ $\to$ \GREEN\ (broken-boost resolution).
\item Agent 2: $r$ suppression \YELLOW\ $\to$ \RED\ (frame confusion; correction is $10^{-38}$).
\item Agent 7: P-Q2 (Born rule) \YELLOW\ $\to$ \GREEN; P-Q3 (einselection) \YELLOW\ $\to$ \GREEN.
\item Agent 3: P-BH3 (GW echoes) \NEW\ \YELLOW.
\end{itemize}

\subsection{Updated Scorecard}

\begin{center}
\begin{tabular}{lccc}
\toprule
\textbf{Metric} & \textbf{i37} & \textbf{i38} & \textbf{Change} \\
\midrule
\GREEN\ predictions & 53 & 56 & $+3$ \\
\YELLOW\ predictions & 7 & 5 & $-2$ \\
\RED\ predictions & 4 & 5 & $+1$ \\
Total predictions & 64 & 66 & $+2$ (1 new, 1 split) \\
\midrule
Positivity bounds & \YELLOW & \GREEN & Resolved \\
Inflation coupling & Unknown & Inert spectator & Characterized \\
$r$ suppression & \YELLOW & \RED & Frame error corrected \\
Born rule & \YELLOW & \GREEN & Quantum corrections negligible \\
Einselection & \YELLOW & \GREEN & Robust \\
GW echoes & --- & \YELLOW & New prediction \\
\bottomrule
\end{tabular}
\end{center}

\subsection{Complete Prediction Table (Abridged)}

\begin{center}
\begin{longtable}{p{1cm}p{6cm}p{1.5cm}p{4cm}}
\toprule
\textbf{ID} & \textbf{Prediction} & \textbf{Grade} & \textbf{Notes} \\
\midrule
\endhead
\multicolumn{4}{l}{\textbf{MOND Sector (22 predictions)}} \\
\midrule
M1 & $\mu(\chi) = \chi/(1+\chi)$ & \GREEN & Derived from $G_2$ \\
M2 & BTFR: $M \propto v^4$ & \GREEN & Parameter-free \\
M3 & RAR: $g_{\rm obs} = g_{\rm bar}\nu(g_{\rm bar}/a_0)$ & \GREEN & 2693 galaxies \\
M4 & EFE existence & \GREEN & Detected at $4\sigma$ \\
M5 & EFE quantitative & \YELLOW & Conflicting data \\
M6--M22 & (remaining galaxy predictions) & 16$\times$\GREEN, 1$\times$\YELLOW & Standard MOND tests \\
\midrule
\multicolumn{4}{l}{\textbf{Quantum Sector (12 predictions)}} \\
\midrule
Q1 & Constraint-field quantization & \GREEN & $A_0$-analogy holds \\
Q2 & Born rule derivation & \GREEN & \textbf{Upgraded i38} \\
Q3 & Position einselection & \GREEN & \textbf{Upgraded i38} \\
Q4 & BMV $2200\times$ enhancement & \GREEN & Testable 2030s \\
Q5 & Lattice Monte Carlo & \YELLOW & Needs implementation \\
Q6 & Zero scalar radiation & \GREEN & Testable with SKA \\
Q7 & $m_g = 0$ (exact) & \GREEN & Shift symmetry \\
Q8--Q12 & (decoherence rates, etc.) & 5$\times$\GREEN & Standard \\
\midrule
\multicolumn{4}{l}{\textbf{Black Hole Sector (10 predictions)}} \\
\midrule
BH1 & Stealth Schwarzschild & \GREEN & Exact solution \\
BH2 & Non-stealth Kerr & \YELLOW & Needs numerics \\
BH3 & GW echoes $D_0 = 0.017$ & \YELLOW & \textbf{New i38}. At LVK boundary \\
BH4 & $T_H^{\rm DFD} = 0.91\,T_H^{\rm GR}$ & \GREEN & From $D_0$ \\
BH5--BH10 & (entropy, no-hair, etc.) & 6$\times$\GREEN & Standard \\
\midrule
\multicolumn{4}{l}{\textbf{Cosmological Sector (12 predictions)}} \\
\midrule
C1 & $n_s$ correction $\sim 10^{-4}$ & \GREEN & Undetectable but consistent \\
C2 & $r$ correction $\sim 10^{-38}$ & \RED & \textbf{Downgraded i38}. Frame error \\
C3 & CMB third peak & \YELLOW & Unsolved \\
C4 & SN Ia consistency & \GREEN & $< 10^{-5}$ deviation \\
C5 & BAO consistency & \GREEN & $< 10^{-5}$ deviation \\
C6--C12 & (various cosmological) & 5$\times$\GREEN, 2$\times$\RED & \\
\midrule
\multicolumn{4}{l}{\textbf{Solar System / GR Tests (10 predictions)}} \\
\midrule
SS1 & $\gamma = 1$ (PPN) & \GREEN & Exact (shift symmetry) \\
SS2 & $\beta = 1$ (PPN) & \GREEN & Exact (shift symmetry) \\
SS3--SS10 & (Shapiro, precession, etc.) & 8$\times$\GREEN & All passed \\
\bottomrule
\end{longtable}
\end{center}

\subsection{The 5 RED Predictions}

\begin{enumerate}
\item \textbf{P-VSL1:} $n_s \approx 0.967$ from VSL. \RED\ (DFD cannot replace inflation).
\item \textbf{P-DE1:} $\psi$-screen explains dark energy. \RED\ ($\psi$-screen $10^4\times$ too small).
\item \textbf{P-DE2:} No bare $\Lambda$ needed. \RED\ (DFD requires $\Lambda$).
\item \textbf{P-INF1:} $r$ suppressed by $\sim 15\%$. \RED\ (frame confusion; actual correction $10^{-38}$).
\item \textbf{P-LSS1:} $\nu$HDM large-scale structure. \RED\ (Katz et al.\ 2026 rules out at $5\sigma$, inherited from MOND).
\end{enumerate}

All 5 RED predictions are in the cosmological sector. DFD's core domain (galaxies, quantum, BHs, solar system) has \textbf{zero RED predictions}.

\subsection{The 5 YELLOW Predictions}

\begin{enumerate}
\item \textbf{P-CMB1:} CMB third peak. Trending \RED.
\item \textbf{P-BH2:} Non-stealth Kerr. Needs numerics.
\item \textbf{P-BH3:} GW echoes from $D_0$. At LVK sensitivity boundary.
\item \textbf{P-MOND5:} EFE quantitative. Conflicting data.
\item \textbf{P-Q5:} Lattice Monte Carlo. Needs implementation.
\end{enumerate}

\end{document}
